This section is devoted to checking that {\hyrql} satisfies usual properties of
typed languages. Moreover, it studies the complexity of checking the
orthogonality predicate. We also prove that we can compile an expressive subset
of terms of \hyrql{} into quantum circuits, with a guarantee on the size of the
circuit. This gives us a characterization of the class of functions computable
in quantum polynomial time, known as \fbqp{}~\cite{BV97}.

\subsection{Standard Properties}
\label{subsec:properties}


By design of the type system, the non-linear context $\Gamma$ can be weakened
and typing still holds. This result does not hold for $\Delta$, as intended,
due to its linear behavior.

\begin{restatable}[Weakening]{proposition}{weakening}
  \label{lem:weakening}
  Let $\typing \Gamma \Delta t T$ be a well-typed term.
  Then $\typing {\Gamma, \Gamma'} \Delta t T$ holds for any context
  $\Gamma'$ compatible with $\Gamma, \Delta$.
\end{restatable}

The reduction relation $\reduces$ is confluent up to equivalence for
well-typed terms.

\begin{restatable}[Confluence]{theorem}{confluence}
  \label{thm:confluence}
  Given a well-typed term $t$, if there exist $t_1$ and $t_2$ such
  that $t \sreduces t_1$ and $t\sreduces t_2$,
  then there exist $t_3$ and $t_4$
  such that $t_1 \sreduces t_3$, $t_2 \sreduces t_4$ and $t_3 \equiv t_4$.
\end{restatable}
Due to rule (Equiv) of Table~\ref{tab:reduction}, confluence is up to
equivalence $\equiv$. Moreover, confluence only holds on well-typed terms.
A consequence of Theorem~\ref{thm:confluence} is that any terminating term has
a unique normal form, up to equivalence. This will be useful for progress
(Theorem~\ref{lem:progress}).

\begin{restatable}[Canonical form for typed terms]{lemma}{canonicaltyped}
  \label{lem:canonical-typed}
  Let $\typing \Gamma \Delta t T$ be a well-typed term.
  Then $t$ has a canonical form $\sum_{i=1}^n \alpha_i \cdot p_i$,
  and this canonical form is unique, up to equivalence on each $p_i$
  and term reordering. Furthermore, $\forall\ 1 \leq i \leq n,\ \typing \Gamma
  \Delta {p_i} T$.
\end{restatable}

As there are no distinct and equivalent pure values,
Lemma~\ref{lem:canonical-typed} implies that the canonical forms of
Definition~\ref{def:orthogonality} exist and are unique.

Typing implies that the values are the normal forms
of the language, up to equivalence.

\begin{restatable}[Progress]{theorem}{progress}
  \label{lem:progress}
  Let $\closedtyping t T$ be a closed term, either $t$ is equivalent
  to a value, or $t$ reduces.
\end{restatable}

Contrary to standard progress lemmas, we require $t$ to be a value up to
equivalence. For instance, the term $\zket + 0 \cdot (\lbd y y) \oket$ does not 
reduce and is not a value.

Finally, typing is also preserved by reduction, provided we consider pure or
terminating terms.

\begin{restatable}[Subject reduction]{theorem}{subred}
  \label{lem:subred}
  Let $\typing \Gamma \Delta t T$ be a well-typed term, and $t \reduces t'$.
  If $t$ terminates or $t$ is pure,
  then $\typing \Gamma \Delta {t'} T$.
\end{restatable}

\begin{remark}\label{rem:subred}
  Theorem~\ref{lem:subred} does not hold on non-terminating and non-pure
  terms, as we cannot type non-terminating superpositions.
  This is the case with the following example, with the Bell state
  $\ket{\Phi^+} =
  \frac{1}{\sqrt 2} \cdot \zket\otimes \zket
  + \frac{1}{\sqrt 2} \cdot \oket\otimes \oket$:
  \[
    \left(\lbd y \qcase{
        ((\letrec f x fx)\zket)
    }{\zket \otimes y}{\oket \otimes y}\right) \ket{\Phi^+}
  \]
  While such a term is well-typed, it reduces to a superposition of
  non-terminating pure terms, which is impossible to type through the rule
  \hyperlink{typesup}{(sup)} from Table~\ref{tab:typing}, as
  the orthogonality predicate requires termination. This is a counterpart of
  having a general language with few restrictions; however, as resource
  analysis will be done on terminating terms, such ill-behaved cases will not be
  considered.
\end{remark}

\subsection{Decidability of Orthogonality}
\label{subsec:orthogonality}
The orthogonality predicate of Definition~\ref{def:orthogonality} is
extensionally complete. As a consequence, it is undecidable.
Indeed, it is as hard as the Universal Halting Problem~\cite{EGZ09}, that is
$\Pi^0_2$-complete in the arithmetical hierarchy~\cite{Nie09}.
Orthogonality uses the inner product, which requires computing
additions, multiplications, and nullity checks on complex numbers. Recall that
algebraic numbers are complex numbers in $\mathbb C$ that are roots of a
polynomial in $\mathbb Q[X]$, and write their sets as $\bar{\mathbb C}$. In the
field of algebraic numbers, equality is decidable, and product and sum are
computable in polynomial time~\cite{HHK05}. We thus restrict ourselves to terms
that only contain amplitudes $\alpha \in \bar{\mathbb C}$.

\begin{restatable}[Undecidability of orthogonality]{theorem}{orthoundecide}
  \label{thm:orthogonality-undecidable}
  Deciding orthogonality between two well-typed terms is $\pitwo$-complete.
\end{restatable}

This is not surprising as \hyrql{} allows for
general recursion. Decidability can be recovered under some slight restrictions.
Towards that end, we introduce for basic types
a \emph{type depth} $d$ as a partial function $d: \{\qbit \}\cup\bset \to
\mathbb N $, defined as follows:
\[
  d(\qbit) \triangleq 1 \qquad
  d(B) \triangleq \max_{\substack{{ s = \sign{\bkvar_1, \dots, \bkvar_n}
  \bvar} \\ c^s \in \cons B}}
  (\sum_{i=1}^n d(\bkvar_i)) + 1
\]
Note that $d(\bkvar)$ is undefined on inductive types.
We call \emph{finite types} any basic type with a defined depth.

\begin{restatable}{proposition}{ortholower}
  \label{prop:orthogonality-lower}
  Given a finite type $\bkvar$, and two terminating terms $\closedtyping s
  \bkvar$, $\closedtyping t \bkvar$, it is decidable whether they are
  orthogonal.
\end{restatable}

In previous works of \spm~\cite{CSV23, CLV24},
the syntax of the language was restrictive enough to ensure the decidability of
orthogonality. As \spm{} is a strict sublanguage of {\hyrql}, orthogonality
can also be decided on this fragment.

\subsection{Complexity Properties}\label{subsec:poly}

We finally discuss how quantum circuits and their size are related to terms of
\hyrql{} and their runtime-complexity. Towards that end, we introduce a
sublanguage of \hyrql{} ($\hyrqlbound T$, Definition~\ref{def:hyrqlbound})
based on three restrictions:
one on the runtime-complexity of the program, to ensure terminating terms
($\timeset T$, Definition~\ref{def:time});
one on types, to obtain terms that can be represented as circuits
($\circuittype$, Definition~\ref{def:restrict-circuit});
and one on recursive calls, to have a faithful relation between the size of the
circuit and the runtime-complexity ($\boundedrec$, Definition~\ref{def:brec}).
This allows, for terms terminating in polynomial time, to characterize
precisely the class \fbqp{}. We first define the time termination of a term.

\begin{definition}\label{def:time}
  Given a term $\closedtyping t Q \linfunc Q'$ and a function $T : \mbN \to \mbN$,
  we say that $t$ \emph{terminates in time $T$}, if for any value 
  $\closedtyping v Q$, $tv$ terminates in at most $T(\size v)$ steps.
  Let $\timeset T$ be the set of terms terminating in time $T$.
\end{definition}

Note that this definition should not be confused with termination defined in
Section~\ref{subsec:semantics}.

To obtain terms that can be compiled to a circuit form, we introduce a 
typed restriction \circuittype{}, which relies on the following type grammar:
\[
    B_r \Coloneqq \qbit \gmid \nat \gmid \typelist{B_r} \gmid B_r \otimes
    B_r \gmid B_r \linfunc B_r \gmid B_r \nonlinfunc B_r
\]

\begin{definition}[Circuit terms]\label{def:restrict-circuit}
  We define $\circuittype$ as the set of well-typed terms $\pi \triangleright
  \closedtyping s Q \linfunc Q'$ satisfying the following conditions:
  \begin{itemize}
    \item Any type in $\pi$ is generated by the grammar of $B_r$;
    \item Any rule $\pi \rightslice (\mathrm{sup})[t_1, \dots, t_n]$
      satisfies $t_i \in \valueset$;
    \item Any rule $\pi \rightslice (\mathrm{qcase})[t, t_0, t_1]$ satisfies
      either $t_0, t_1 \in \valueset$; or $t_0 = c(\ket 0, s_0)$ and
      $t_1 = c(\ket 1, s_1)$, for $c \in \set{\otimes, ::}$.
  \end{itemize}
\end{definition}

While the first condition yields terms that can be represented
by a circuit, the two others allow for an efficient compilation and 
guarantees that any superposition can be compiled effectively to circuits.
Note that while these terms only take inputs of type $Q$, to represent
them clearly in a quantum circuit, they can nonetheless contain
multiple inputs and higher-order constructs in their body.
This restriction alone is sufficient to compile $t$ to a family of circuits in
Theorem~\ref{thm:circuit-exist}, provided $t$ terminates for a given time $T$.

\begin{example}
  \label{ex:bqwalk-adapt}
  The following term, which is an adaptation of the program from
  Example~\ref{ex:bqwalk}, belongs to $\circuittype$. 
  \[
    \pi \triangleright \closedtyping {\lbd x \match x {q \otimes n \to
    \bqwalk\,q\,n}} {\qbit \otimes \nat \linfunc \typelist \qbit}
  \]
  Any type in $\pi$ is derived from $B_r$. Term superposition happens only in
  $\mathtt{Had}$ and is between values. Furthermore, there are two qcase
  constructs, either between values, or between two terms $\zket :: s_0$ and
  $\oket :: s_1$, thus it satisfies all properties of $\circuittype$. Note that
  as $B_r$ allows for tensor products, any program taking multiple inputs can
  always be rewritten using the same technique so that it belongs to
  $\circuittype$.
\end{example}

In the literature, family of circuits are parameterized by the size
of their input. However, in \hyrql{}, knowing the size might not be sufficient,
e.g., on inputs of type $\typelist \qbit \otimes \typelist \qbit$. 
To remedy this, we will parameterize the family by the \emph{shape} of the
input, which will give us the necessary information. Formally, we define
\[
  \begin{aligned}
    \shapeset Q &\triangleq \set{w \mid \closedtyping w \shape Q} \\
    \shapeset w &\triangleq \set{v \mid \closedtyping v Q \wedge \shape v
    \sreduces w}, \forall w \in \shapeset Q,
  \end{aligned}
\]
where $\shapeset Q$ is the \emph{set of shape of terms of type $Q$}, and
$\shapeset w$ the \emph{set of values of shape $w$}.

Assuming an encoding of $t \in \circuittype$ to a quantum state $\ket t$, which
can be done naturally, a circuit $C$ is said to \emph{approximate} $v$ on input
$v'$ with probability $p$, if $C$ evaluates to $\ket \phi$, on input $\ket{v'}$
(up to ancillary qubits), and $\lvert \langle v | \phi \rangle \rvert^2 \geq
p$.

\begin{restatable}[Circuit compilation]{theorem}{circuitexistence}
  \label{thm:circuit-exist}
  Let $t \in \circuittype \cap \timeset T$ be a term of type $Q \linfunc Q'$
  for a given $T : \mbN \to \mbN$. 
  Then, we can generate a family of circuits $\famcirc \triangleq (C_w)_{w \in
  \shapeset Q}$ such that for any $w \in \shapeset Q$ and any $v \in \shapeset
  w$, $C_w$ approximates $v'$ with probability $\frac 2 3$ on input $v$, where
  $tv \sreduces v'$.
\end{restatable}

We are now interested in the \emph{size} of each circuit of the family, i.e.,
the total number of gates acting on the circuit; in particular, the goal is to
exhibit a relation between the size and the runtime of the initial term $t$. To
avoid a blowup due to recursive calls, we introduce a syntactic restriction
$\boundedrec$. It relies on the notion of the \emph{width} of a variable $f$ in
a term $t$, written $w(f, t)$ and defined inductively as follows:
\[
\begin{array}{c}
  w(f, x) \triangleq \delta_{x, f} 
  \qquad 
  w(f, \ket i) \triangleq 0 
  \qquad 
  w(f, \lbd x t) \triangleq w(f,t) \qquad
  w(f, \letrec g x t) \triangleq w(f, t) \qquad \\[1ex]
  w(f, t_1 t_2) \triangleq w(f, t_1) + w(f, t_2)
  \qquad 
  w(f, \sum_{i=1}^n \alpha_i \cdot t_i) = \max_{1 \leq i \leq n} w(f, t_i) \\[1ex]
  w(f, c(t_1, \dots, t_n))       \triangleq \sum_{i=1}^n w(f, t_i) 
  \qquad
  w(f, \shape t) = w(f, t) \\[1ex]
  w(f,\qcasedefault t) \triangleq w(f, t) + \max(w(f,t_0), w(f,t_1)) \\[1ex]
  w(f, \matchdefaultcollapsed t) \triangleq w(f, t) + \max_{1 \leq i \leq
  n} (w(f, t_i))
\end{array}
\]

\begin{definition}\label{def:brec}
  A term $t$ is said to \emph{compile faithfully}, if for any term $t'=(\letrec
  f x s)$ such that $t' \subtermeq t$, the following conditions hold:
  \begin{itemize}
    \item $w(f, s) \leq 1$;
    \item $\exists !\  t', \forall\ t_1, f t_1 \subtermeq s \implies t_1 = t'$;
    \item $\neg \exists\  t', t'f \subtermeq s$.
  \end{itemize}
  Let $\boundedrec$ be the set of terms that compile faithfully.
\end{definition}

The first two conditions are derived from \cite{HPS25}, and ensure that, for
each recursive program $\letrec f x s$, all recursive calls are done in
different branches of a \textqcase{} or \textmatch{}, with the same input.
Furthermore, the last condition ensures that this behaviour is preserved, even
in presence of higher-order constructs.

\begin{example}
  The program of Example~\ref{ex:bqwalk-adapt} belongs to
  $\boundedrec$, as both $\bqwalk$ and $\tt repeat$ only perform one recursive
  call. All the examples across the paper also satisfy this criterion.
\end{example}

\begin{definition}\label{def:hyrqlbound}
  Let $T : \mbN \to \mbN$. We define $\hyrqlbound T$ as follows:
  \[
    \hyrqlbound T\ \triangleq\ \circuittype\ \cap\ \boundedrec\ \cap\ \timeset T
  \]
\end{definition}

Any term in this subset compiles to quantum circuits of size bounded
polynomially by $T$.

\begin{restatable}[Bounded circuit compilation]{theorem}{circuitbound}
  \label{thm:circuit-bound}
  Let $T : \mbN \to \mbN$, and let $t \in \hyrqlbound T$ be a term of type
  $Q \linfunc Q'$. Then, we can generate a family of circuits
  $\famcirc \triangleq (C_w)_{w \in \shapeset Q}$ and a polynomial $P$ such
  that, for any $w \in \shapeset Q$:
  \begin{itemize}
    \item $\size{C_w} \leq P(T(\size w))$;
    \item $\forall\ v \in \shapeset w$, $C_w$ approximates $v'$ with
      probability $\frac 2 3$ on input $v$, where $tv \sreduces v'$.
  \end{itemize}
\end{restatable}

We now want to use Theorem~\ref{thm:circuit-bound} to obtain a characterization
of \fbqp{} (Definition~\ref{def:fbqp}). We recall that a family of circuits
$(C_n)_{n \in \mbN}$ is said to be \emph{uniform polynomially-sized} with
$\alpha : \mbN \to \mbN$ if there exists $P \in \mbN[X]$ such that $\size{C_n}
\leq P(n)$ and $C_n$ has exactly $\alpha(n)$ ancillary qubits, and there is a
polynomial-time Turing machine, which takes $n$ as input, and outputs a
representation of $C_n$ for any $n \in \mbN$.

\begin{definition}[\cite{BV97}]
  \label{def:fbqp}
  A binary function $f : \binfunc$ is said to be computed by a family of
  circuits $(C_n)_{n \in \mbN}$ if for any $x \in \set{0,1}^*$, $C_{\size x}$
  approximates $f(x)$ with probability $\frac 2 3$ on input $x$. \fbqp{} is
  defined as the set of binary functions computed by a uniform
  polynomially-sized family of circuits.
\end{definition}

To characterize \fbqp{}, we define the following set of terms, which terminate 
in polynomial time, and have list of qubits as input:
\[
  \hyrqlpoly \triangleq \cup_{P \in \mbN[X]} \set{
    t \in \hyrqlbound P \gmid \closedtyping t \typelist \qbit \linfunc Q
  }
\]

These generated families of circuits will be indexed by natural numbers
instead of shape, as one can remark that $\shapeset{\typelist \qbit}$ is in
bijection with $\mbN$.
Altogether, $\hyrqlpoly$ contain terms that are
compiled to circuits approximating functions of $\fbqp$.

\begin{restatable}[\fbqp{} soundness]{theorem}{fbqpsound}
  \label{thm:fbqp-sound}
  Let $t \in \hyrqlpoly$, and let $f : \binfunc$. If $f$ is computed by
  $\famcirc$, then $f \in \fbqp$.
\end{restatable}

For example, the Quantum Fourier Transform program from Example~\ref{ex:qft}
belongs to $\fbqp$, and can thus be compiled to circuits of polynomial size.

We now prove \fbqp{} completeness, meaning that every function of \fbqp{} can
be represented by a function in $\hyrqlpoly$. This is proven by writing
any function of Yamakami's algebra~\cite{Yam20}, which is known to capture
\fbqp{}, into $\hyrqlpoly$.

\begin{restatable}[\fbqp{} completeness]{theorem}{fbqpcomplete}
  \label{thm:fbqp-complete}
  Let $f : \binfunc \in \fbqp$. Then, there exists $t \in \hyrqlpoly$
  such that $\famcirc$ computes $f$.
\end{restatable}

We have given a typed, syntactic and semantic restriction of \hyrql{}, that
compiles into a bounded family of quantum circuits. Such restriction could be
refined, either to accept more terms by lifting requirements in
Definition~\ref{def:brec}, as done in~\cite{HPS25}; or by adding requirements
to characterize more precise complexity classes~\cite{FHPS25}. We are
nonetheless able to encapsulate \fbqp{}, which demonstrates the expressivity of
the $\hyrqlpoly$ fragment.
